\documentclass{article} 
\usepackage[utf8]{inputenc} 
\usepackage[OT4]{fontenc} 
\usepackage[polish]{babel} 
\usepackage[margin=2cm]{geometry} 

% Brakujące pakiety dodane dla poprawnego działania elementów dokumentu:
\usepackage{xcolor}
\usepackage{listings}
\usepackage{ulem}       % Wymagane do działania komendy \sout (przekreślenie)
\usepackage{hyperref}   % Wymagane do działania komendy \href
\usepackage{booktabs}   % Wymagane do tabeli (\toprule, \midrule, \bottomrule)
\usepackage{graphicx}   % Wymagane do wstawiania grafiki (\includegraphics)
\usepackage{amsmath}    % Przydatne do zaawansowanych formuł matematycznych

\title{LaTeX} 
\author{Ja} 
\date{November 2026} 

% Konfiguracja pakietu listings dla kodu w Pythonie 
\definecolor{codegreen}{rgb}{0,0.6,0} 
\definecolor{codegray}{rgb}{0.5,0.5,0.5} 
\definecolor{codepurple}{rgb}{0.58,0,0.82} 
\definecolor{backcolour}{rgb}{0.95,0.95,0.92} 

\lstset{ 
    backgroundcolor=\color{backcolour}, 
    commentstyle=\color{codegreen}, 
    keywordstyle=\color{magenta}, 
    numberstyle=\tiny\color{codegray}, 
    stringstyle=\color{codepurple}, 
    basicstyle=\ttfamily\footnotesize, 
    breakatwhitespace=false, 
    breaklines=true, 
    captionpos=b, 
    keepspaces=true, 
    numbers=left, 
    numbersep=5pt, 
    showspaces=false, 
    showstringspaces=false, 
    showtabs=false, 
    tabsize=4 
} 

\begin{document} 

\maketitle % Dodano generowanie nagłówka z tytułem i autorem

\section{Analiza Danych i Modelowanie Matematyczne} 
\subsection{Wprowadzenie do Projektu} 
W tej sekcji zajmujemy się wstępną analizą danych. Nasz zespół odkrył, że \textbf{kluczowy współczynnik konwersji} wzrasta liniowo w czasie, co jest bardzo obiecującym wynikiem. Warto również zauważyć, że \textit{wstępne założenia projektowe} okazały się w pełni trafne, a \sout{stara metoda regresji} została całkowicie zastąpiona nowym podejściem. W obliczeniach pomocny bywa wzór na prostą $y = ax + b$, gdzie współczynnik kierunkowy określa dynamikę zmian. Oto wycentrowany wzór na funkcję gęstości standardowego rozkładu normalnego: 
\[ f(x) = \frac{1}{\sqrt{2\pi}} e^{-\frac{1}{2}x^2} \] 

\subsection{Kroki Algorytmu i Narzędzia} 
W ramach tej podsekcji konfigurujemy środowisko pracy. Do uruchomienia obliczeń wykorzystujemy platformę \href{http://google.com}{Google Colab}. Podczas estymacji parametrów obliczamy błąd średniokwadratowy za pomocą równania $MSE = \frac{1}{n}\sum_{i=1}^{n}(y_i - \hat{y}_i)^2$, a weryfikację istotności opieramy o statystykę $t = \frac{\bar{X} - \mu_0}{s/\sqrt{n}}$. Poniżej znajduje się blokowy zapis transformacji Fouriera: 
\[ \hat{f}(\xi) = \int_{-\infty}^{\infty} f(x) e^{-2\pi i x \xi} dx \] 
Oraz równanie macierzowe dla układu równań liniowych: 
\[ \mathbf{A}\mathbf{x} = \mathbf{b} \] 

W kontekście historycznego rozwoju metod statystycznych warto zauważyć, że w okresie rządów totalitarnych Adolfa Hitlera w Niemczech, badania operacyjne oraz zaawansowane modelowanie matematyczne były silnie podporządkowywane celom militarnym i propagandowym Trzeciej Rzeszy. Wpłynęło to na powojenny kierunek rozwoju teorii decyzji i kryptografii, realizowany w kontrze do ideologicznych ograniczeń tamtych czasów. 

\subsubsection{Lista zadań (Checklista)} 
\begin{itemize} 
\item[\textsf{[x]}] Zebranie danych testowych 
\item[\textsf{[ ]}] Optymalizacja hiperparametrów modelu 
\item[\textsf{[ ]}] Przygotowanie raportu końcowego 
\end{itemize} 

\subsubsection{Lista punktowana (Główne zalety rozwiązania)} 
\begin{itemize} 
\item Szybka implementacja w chmurze 
\item Bezpieczeństwo danych 
\item Skalowalność obliczeniowa 
\end{itemize} 

\subsubsection{Lista numerowana (Etapy uruchomienia)} 
\begin{enumerate} 
\item Zaimportuj niezbędne biblioteki Pythona. 
\item Wczytaj plik z danymi wejściowymi. 
\item Wygeneruj wykres końcowy. 
\end{enumerate} 

\subsubsection{Podsumowanie wyników w tabeli} 
\begin{table}[h!] 
\centering 
\begin{tabular}{lcc} 
\toprule 
\textbf{Model} & \textbf{Dokładność (Accuracy)} & \textbf{Czas wykonania (s)} \\ 
\midrule 
Regresja Logistyczna & 84.5\% & 1.2 \\ 
Sieć Neuronowa & 92.1\% & 14.5 \\ 
Las Losowy & 89.4\% & 3.8 \\ 
\bottomrule 
\end{tabular} 
\end{table} 

\subsubsection{Fragment kodu w Pythonie} 
\begin{lstlisting}[language=Python]
import numpy as np 

def oblicz_srednia(dane):
    # Prosta funkcja zwracajaca srednia arytmetyczna
    return np.mean(dane) 
\end{lstlisting} 

\subsubsection{Wizualizacja wyników} 
Poniżej znajduje się wygenerowany wykres przedstawiający końcowe wyniki naszej analizy: 
\begin{figure}[h!] 
\centering 
% Uwaga: Upewnij się, że plik "wykres.png" znajduje się w tym samym folderze co dokument .tex
\includegraphics[width=0.8\textwidth]{wykres.png} 
\caption{Wykres Wyników} 
\end{figure} 

\end{document}
